\subsection{点处的可微函数}

2.1.1. 设 \(f\) 在 E 中一点 \(x_{0}\) 的某邻域上定义且取值于 F。设 \(u\) 是从 E 到 F 的连续线性映射空间 \(\mathcal{L}(\mathrm{E};\mathrm{F})\) 中的元素。\(f\) 在 \(x_{0}\) 处可微且在那里以 \(u\) 为导数的充要条件是有

\[
\lim _ {h \rightarrow 0, h \neq 0} \frac {f (x _ {0} + h) - f (x _ {0}) - u (h)}{\| h \|} = 0.
\]

2.1.2. \(f\) 在 \(x_{0}\) 处严格可微的充要条件是有

\[
\lim_{\substack{(h,k)\to (0,0)\\ h\neq k}}\frac{f(x_{0} + h) - f(x_{0} + k) - \mathrm{D}f(x_{0})\cdot(h - k)}{\|h - k\|} = 0.
\]

2.1.3. 设 \(F_{1}\) 和 \(F_{2}\) 是两个分离局部凸空间，u 是从 \(F_{1} \times F_{2}\) 到 F 的双线性映射，并满足下列连续性条件：

(SC) 若 \(((a_{n}, b_{n}))\) 是 \(F_{1} \times F_{2}\) 中收敛于元素 \((a, b) \in \mathbf{F}_{1} \times \mathbf{F}_{2}\) 的元素序列，则序列 \((u(a_{n}, b_{n}))\) 在 F 中收敛于 \(u(a, b)\)。

设 \(f_{i}\)（\(i = 1, 2\)）是从 E 中一点 \(x_{0}\) 的某邻域到 \(F_{i}\) 的映射。若 \(f_{1}\) 和 \(f_{2}\) 在 \(x_{0}\) 处可微，则 \(u(f_{1}, f_{2})\) 在 \(x_{0}\) 处可微，且有：

\[
\mathbf {D} (u (f _ {1}, f _ {2})) (x _ {0}). h = u (\mathbf {D} f _ {1} (x _ {0}). h, f _ {2} (x _ {0})) + u (f _ {1} (x _ {0}), \mathbf {D} f _ {2} (x _ {0}). h)
\]

对所有 \(h \in E\) 均成立。

